% Options for packages loaded elsewhere
\PassOptionsToPackage{unicode}{hyperref}
\PassOptionsToPackage{hyphens}{url}
\PassOptionsToPackage{dvipsnames,svgnames,x11names}{xcolor}
%
\documentclass[
  10pt,
]{article}
\usepackage{amsmath,amssymb}
\usepackage{iftex}
\ifPDFTeX
  \usepackage[T1]{fontenc}
  \usepackage[utf8]{inputenc}
  \usepackage{textcomp} % provide euro and other symbols
\else % if luatex or xetex
  \usepackage{unicode-math} % this also loads fontspec
  \defaultfontfeatures{Scale=MatchLowercase}
  \defaultfontfeatures[\rmfamily]{Ligatures=TeX,Scale=1}
\fi
\usepackage{lmodern}
\ifPDFTeX\else
  % xetex/luatex font selection
  \setmainfont[]{DejaVu Serif}
  \setsansfont[]{DejaVu Sans}
  \setmonofont[]{DejaVu Sans Mono}
  \setmathfont[]{Latin Modern Math}
\fi
% Use upquote if available, for straight quotes in verbatim environments
\IfFileExists{upquote.sty}{\usepackage{upquote}}{}
\IfFileExists{microtype.sty}{% use microtype if available
  \usepackage[]{microtype}
  \UseMicrotypeSet[protrusion]{basicmath} % disable protrusion for tt fonts
}{}
\makeatletter
\@ifundefined{KOMAClassName}{% if non-KOMA class
  \IfFileExists{parskip.sty}{%
    \usepackage{parskip}
  }{% else
    \setlength{\parindent}{0pt}
    \setlength{\parskip}{6pt plus 2pt minus 1pt}}
}{% if KOMA class
  \KOMAoptions{parskip=half}}
\makeatother
\usepackage{xcolor}
\usepackage[margin=2.0cm]{geometry}
\usepackage{longtable,booktabs,array}
\usepackage{calc} % for calculating minipage widths
% Correct order of tables after \paragraph or \subparagraph
\usepackage{etoolbox}
\makeatletter
\patchcmd\longtable{\par}{\if@noskipsec\mbox{}\fi\par}{}{}
\makeatother
% Allow footnotes in longtable head/foot
\IfFileExists{footnotehyper.sty}{\usepackage{footnotehyper}}{\usepackage{footnote}}
\makesavenoteenv{longtable}
\setlength{\emergencystretch}{3em} % prevent overfull lines
\providecommand{\tightlist}{%
  \setlength{\itemsep}{0pt}\setlength{\parskip}{0pt}}
\setcounter{secnumdepth}{-\maxdimen} % remove section numbering
% doc_header.tex -- LaTeX preamble additions used by build_docs.sh for every
% Markdown -> LaTeX -> PDF conversion (supplements, README, submission files).
\usepackage{etoolbox}
% longtable may not be loaded by pandoc when a document has no table
\usepackage{longtable}
% bold vectors: \bm -> \symbf under unicode-math (XeLaTeX), bm package otherwise
\makeatletter
\@ifpackageloaded{unicode-math}{\providecommand{\bm}[1]{\symbf{#1}}}{\usepackage{bm}}
\makeatother
% smaller type in long tables so that wide audit tables fit the text width;
% inside a table a line break is allowed after an underscore of plain text
% (pandoc writes it as \_), so that long identifiers wrap inside their column
% rows are set a little apart, so that symbols with indices do not touch
\AtBeginEnvironment{longtable}{\footnotesize\renewcommand{\arraystretch}{1.45}%
  \let\sptunderscore\_\def\_{\sptunderscore\allowbreak}}
\setlength{\LTcapwidth}{\textwidth}
% allow line breaks after slashes and hyphens inside table cells
\usepackage{xurl}
% code blocks: break over-long lines instead of running into the margin.
% Plain ``` blocks are typeset by pandoc as {verbatim}, highlighted ones
% (```bash) as {Highlighting}; both are re-declared as breakable fancyvrb
% environments.  The continuation line is marked with a small hooked arrow.
\usepackage{fvextra}
\RecustomVerbatimEnvironment{verbatim}{Verbatim}{breaklines,breakanywhere,fontsize=\small}
\makeatletter
\@ifundefined{Highlighting}{}{%
  \RecustomVerbatimEnvironment{Highlighting}{Verbatim}{breaklines,breakanywhere,fontsize=\small,commandchars=\\\{\}}}
\makeatother
% inline code (pandoc writes it as \texttt{...}): never hyphenate -- an inserted
% hyphen could be read as part of a file name -- but allow a line break after
% an underscore, slash, comma, colon or equals sign, so that long paths and
% column lists wrap inside the text width and inside table columns.
% Inside a formula (e.g. the name of a notebook symbol in an equation) the same
% command must first leave math mode, which \text of amsmath does.
\ExplSyntaxOn
\tl_new:N \l__spt_code_tl
\cs_new_protected:Npn \__spt_texttt:n #1
  {
    \tl_set:Nn \l__spt_code_tl {#1}
    \regex_replace_all:nnN { ([\/\,\:\=]) } { \1 \c{allowbreak} } \l__spt_code_tl
    \regex_replace_all:nnN { \c{_} } { \c{_} \c{allowbreak} } \l__spt_code_tl
    \group_begin: \ttfamily \hyphenchar \font = -1 \scan_stop: \l__spt_code_tl \group_end:
  }
\RenewDocumentCommand \texttt { m }
  {
    \mode_if_math:TF { \text { \__spt_texttt:n {#1} } } { \__spt_texttt:n {#1} }
  }
\ExplSyntaxOff
% A figure is written in the Markdown source as an image paragraph followed by
% a caption paragraph ("**Figure S1.** ..."), so that the caption is also
% visible on GitHub.  An image that starts a paragraph is centered.
\usepackage{graphicx}% pandoc loads it only for documents that contain an image
\NewCommandCopy\sptincludegraphics\includegraphics
\RenewDocumentCommand\includegraphics{O{} m}{%
  \ifvmode\begin{center}\sptincludegraphics[#1]{#2}\end{center}%
  \else\sptincludegraphics[#1]{#2}\fi}
% keep figures where they are referenced
\usepackage{float}
\floatplacement{figure}{H}
% running header with the document title
\usepackage{fancyhdr}
\pagestyle{fancy}
\fancyhf{}
\fancyhead[L]{\small PROJECT spacetimes --- documentation and provenance}
\fancyhead[R]{\small\thepage}
\renewcommand{\headrulewidth}{0.4pt}
% 16x16 matrices: tighter columns so that a full matrix fits the text width
\setlength{\arraycolsep}{3.2pt}
% The documents of this project contain long tagged symbols, file paths and
% identifiers that cannot be hyphenated.  A paragraph is therefore allowed to
% have looser lines before TeX gives up and lets a line stick out of the margin.
\AtBeginDocument{\tolerance=2000 \emergencystretch=5em \hbadness=2000}
\ifLuaTeX
  \usepackage{selnolig}  % disable illegal ligatures
\fi
\IfFileExists{bookmark.sty}{\usepackage{bookmark}}{\usepackage{hyperref}}
\IfFileExists{xurl.sty}{\usepackage{xurl}}{} % add URL line breaks if available
\urlstyle{same}
\hypersetup{
  pdftitle={P8 --- The source code: what computes what, and how to compile and run it},
  pdfauthor={Prepared for Patrick Lee Nash, Ph.D.},
  colorlinks=true,
  linkcolor={blue},
  filecolor={Maroon},
  citecolor={blue},
  urlcolor={blue},
  pdfcreator={LaTeX via pandoc}}

\title{P8 --- The source code: what computes what, and how to compile
and run it}
\usepackage{etoolbox}
\makeatletter
\providecommand{\subtitle}[1]{% add subtitle to \maketitle
  \apptocmd{\@title}{\par {\large #1 \par}}{}{}
}
\makeatother
\subtitle{Provenance required by Hard Rule 6 of the task specification
(PROJECT spacetimes)}
\author{Prepared for Patrick Lee Nash, Ph.D.}
\date{November 12, 2026}

\begin{document}
\maketitle

\hypertarget{what-this-document-is-for}{%
\section{1. What this document is for}\label{what-this-document-is-for}}

Hard Rule 6 of the task specification demands commented source code for
every equation, result, figure, graph and prediction, instructions for
using the code, a provenance for every computed value, and that the code
be executed. This document is the map: it says which program computes
which file, which file feeds which figure and which number of the
documents, how the programs are tested, and how to compile and run
everything.

The rules that the code follows:

\begin{itemize}
\tightlist
\item
  \textbf{Rust for every computation.} The task specification prefers
  Rust. All physics is computed by one Rust program, \texttt{spacetimes}
  (\texttt{code/spacetimes-rs}, 28 source files, about 11 400 lines, 63
  unit tests). It uses no external library, so it builds with nothing
  but a Rust compiler.
\item
  \textbf{Python for an independent verification, and for tables,
  figures and documents.} A program can be wrong in the same way twice;
  the Dirac matrices are therefore verified a second time by a Python
  program that shares no code with the Rust program. Python scripts also
  turn the data files into tables and figures and build the documents.
\item
  \textbf{Mathematica is not used.} The matrices are read from what
  Mathematica printed in the author's notebooks (provenance P3) and
  rebuilt by the programs.
\item
  \textbf{No number is typed by hand.} Every number in a document is
  inserted by a script from a data file that a program wrote (Sec. 6).
\end{itemize}

\textbf{Where the code was executed.} All data files, figures and
documents of the repository were produced by running the code on a Linux
computer (x86-64) with rustc 1.97.0, Python 3.13.18 (numpy 2.5.3, scipy
1.18.1, matplotlib 3.11.2), pandoc 3.1.3 and XeLaTeX (TeX Live 2023).
The assistant could not run the Rust program on the author's Mac,
because no Rust compiler is reachable from its sandbox there; since the
program has no dependencies, it should build and run on macOS unchanged
(Sec. 2). Results of floating-point arithmetic may then differ in the
last digits; the exact checks of the matrices use integer arithmetic and
must agree exactly.

\hypertarget{how-to-compile-and-run-everything-on-a-mac}{%
\section{2. How to compile and run everything on a
Mac}\label{how-to-compile-and-run-everything-on-a-mac}}

Each block can be pasted into the Terminal as it is. Lines that begin
with \texttt{\#} are explanations.

\textbf{Step 1, once: the Rust compiler.} If the command
\texttt{cargo\ -\/-version} prints a version number, skip this step.

\begin{verbatim}
# install the Rust compiler and its build tool "cargo" (official installer; accept the defaults)
curl --proto '=https' --tlsv1.2 -sSf https://sh.rustup.rs | sh
# make the new commands known to this Terminal window
source "$HOME/.cargo/env"
cargo --version
\end{verbatim}

\textbf{Step 2, once: the Python packages} (needed for the tables, the
figures and the comparison with the reference spectra, not for the Rust
program).

\begin{verbatim}
python3 -m pip install --user numpy scipy matplotlib
\end{verbatim}

\textbf{Step 3: compile and test.} This compiles the program (about one
minute) and runs its 63 unit tests; the last line must say
\texttt{test\ result:\ ok.\ 63\ passed}.

\begin{verbatim}
cd "/Users/nsh/Developer/github/spacetimes/code/spacetimes-rs"
cargo build --release
cargo test --release
\end{verbatim}

\textbf{Step 4: recompute every data file.} Each command writes its
files into \texttt{data/} and prints what it did. The commands
\texttt{forecast} (about one hour) and \texttt{superfluid} (about ten
minutes) are the long ones; the whole block takes about an hour and a
half on a computer with two processor cores.

\begin{verbatim}
cd "/Users/nsh/Developer/github/spacetimes"
B=code/spacetimes-rs/target/release/spacetimes
for c in matrices channels gap dimer symmetry cjt background spectrum forecast superfluid halos
do
  $B $c
done
\end{verbatim}

\textbf{Step 5: the independent verification and the comparison with
CLASS.}

\begin{verbatim}
cd "/Users/nsh/Developer/github/spacetimes"
python3 code/scripts/verify_matrices_independent.py
python3 code/scripts/validate_cmb_against_class.py
python3 code/scripts/check_symbols.py
\end{verbatim}

\textbf{Step 6: tables and figures.}

\begin{verbatim}
cd "/Users/nsh/Developer/github/spacetimes"
python3 code/scripts/make_tables.py
python3 code/scripts/make_figures.py
\end{verbatim}

\textbf{Step 7, optional: the documents.} This needs pandoc and a TeX
installation with XeLaTeX and latexmk (for example MacTeX), and the
DejaVu and Latin Modern Math fonts.

\begin{verbatim}
cd "/Users/nsh/Developer/github/spacetimes"
for f in docs_src/*.src.md; do python3 code/scripts/assemble_docs.py "$f"; done
bash code/scripts/build_docs.sh
python3 code/scripts/build_manuscripts.py
\end{verbatim}

To run a single command of the program and see the list of all commands:

\begin{verbatim}
cd "/Users/nsh/Developer/github/spacetimes"
code/spacetimes-rs/target/release/spacetimes help
\end{verbatim}

\hypertarget{the-commands-of-the-rust-program-and-the-files-they-write}{%
\section{3. The commands of the Rust program, and the files they
write}\label{the-commands-of-the-rust-program-and-the-files-they-write}}

Every data file begins with lines that start with \texttt{\#} and state
the program, the command, the content, the units and the parameters; the
table can be checked against these headers.

\begin{longtable}[]{@{}
  >{\raggedright\arraybackslash}p{(\columnwidth - 6\tabcolsep) * \real{0.1188}}
  >{\raggedright\arraybackslash}p{(\columnwidth - 6\tabcolsep) * \real{0.4257}}
  >{\raggedright\arraybackslash}p{(\columnwidth - 6\tabcolsep) * \real{0.3168}}
  >{\raggedright\arraybackslash}p{(\columnwidth - 6\tabcolsep) * \real{0.1386}}@{}}
\toprule\noalign{}
\begin{minipage}[b]{\linewidth}\raggedright
Command
\end{minipage} & \begin{minipage}[b]{\linewidth}\raggedright
What it computes
\end{minipage} & \begin{minipage}[b]{\linewidth}\raggedright
Files in \texttt{data/}
\end{minipage} & \begin{minipage}[b]{\linewidth}\raggedright
Documented in
\end{minipage} \\
\midrule\noalign{}
\endhead
\bottomrule\noalign{}
\endlastfoot
\texttt{matrices} & builds the two sets of eight real \(16\times16\)
Dirac matrices and runs the 106 machine checks; classifies the fermion
pairs & \texttt{gamma1644.txt}, \texttt{gamma1671.txt},
\texttt{matrix\_checks.txt}, \texttt{matrix\_checks.csv},
\texttt{pair\_classification.txt}, \texttt{pair\_classification.csv};
also \texttt{provenance/generated/matrices.md} & ``dirac matrices''; M
II to IV \\
\texttt{channels} & exchange amplitudes of the Vorton-Scalar in the pair
channels; quartic invariant of the order parameter &
\texttt{pair\_channel\_kernel.csv}, \texttt{gl\_quartic\_check.csv} & S1
Secs. 4, 6; M VI, VIII \\
\texttt{gap} & gap equation of the hidden-axis singlet &
\texttt{gap\_scan.csv}, \texttt{gap\_profile.csv} & S1 Sec. 5; M VII \\
\texttt{dimer} & two fermions in the Yukawa potential: scattering
length, binding energy, critical coupling & \texttt{dimer.csv} & S1 Sec.
5.3; M VII \\
\texttt{symmetry} & internal symmetry of the zero-mode sector, from the
matrices & \texttt{reduced\_symmetry.csv} & S1 Sec. 2.1; S4 Sec. 2.4; M
V \\
\texttt{cjt} & everything computed with the effective action of
Cornwall, Jackiw and Tomboulis: test model, paired state,
thermodynamics, order-parameter landscape, matter, temperature, vacuum &
the sixteen files \texttt{cjt\_*.csv} & S4; M VII to IX, App. B \\
\texttt{background} & expansion history and recombination with the dark
radiation of the model & \texttt{step\_thermo.csv} & S3 Sec. 1; M XI \\
\texttt{spectrum} & spectra of the microwave background and of matter
for the fiducial model; convergence & \texttt{cmb\_fiducial\_cl.csv},
\texttt{cmb\_fiducial\_pk.csv}, \texttt{cmb\_convergence.csv} & S3 Sec.
2; M XI, App. D \\
\texttt{forecast} & fits of seven families of dark-radiation models to
mock data & \texttt{dr\_fisher.csv}, \texttt{dr\_profiles.csv},
\texttt{dr\_residuals.csv}, \texttt{dr\_zt\_scan.csv},
\texttt{dr\_chord\_check.csv} & S3 Sec. 3; M XI \\
\texttt{superfluid} & growth of structure if the dark matter is a
superfluid & \texttt{sfdm\_transfer.csv}, \texttt{sfdm\_summary.csv},
\texttt{sfdm\_bounds.csv} & S3 Sec. 4; M XII \\
\texttt{halos} & conditions for degeneracy and condensation in halos;
self-interaction of bound pairs & \texttt{dm\_halo\_estimates.csv},
\texttt{dm\_selfinteraction.csv} & S3 Sec. 4; M XII \\
\end{longtable}

Files in \texttt{data/} that other programs write:
\texttt{matrix\_checks\_independent.txt}
(\texttt{verify\_matrices\_independent.py});
\texttt{cmb\_validation.csv}, \texttt{pk\_validation.csv}
(\texttt{validate\_cmb\_against\_class.py});
\texttt{notebook\_matrices\_4+4.txt},
\texttt{notebook\_matrices\_7+1.txt}
(\texttt{tools/extract\_nb\_matrices.py}, provenance P3);
\texttt{symbol\_checks.txt}, \texttt{symbol\_concordance.csv}
(provenance P1); and the two reference spectra in
\texttt{data/reference/}, which are copies from the author's project
NCSS with their checksums in the headers.

\hypertarget{the-source-files-of-the-rust-program}{%
\section{4. The source files of the Rust
program}\label{the-source-files-of-the-rust-program}}

Each file begins with a comment that states its purpose, its equations
and its tests. Four small files are taken unchanged from the author's
project NCSS and say so in their first lines.

\begin{longtable}[]{@{}
  >{\raggedright\arraybackslash}p{(\columnwidth - 4\tabcolsep) * \real{0.2000}}
  >{\raggedleft\arraybackslash}p{(\columnwidth - 4\tabcolsep) * \real{0.0750}}
  >{\raggedright\arraybackslash}p{(\columnwidth - 4\tabcolsep) * \real{0.7250}}@{}}
\toprule\noalign{}
\begin{minipage}[b]{\linewidth}\raggedright
File
\end{minipage} & \begin{minipage}[b]{\linewidth}\raggedleft
Lines
\end{minipage} & \begin{minipage}[b]{\linewidth}\raggedright
Purpose
\end{minipage} \\
\midrule\noalign{}
\endhead
\bottomrule\noalign{}
\endlastfoot
\texttt{gamma.rs} & 270 & construction of the two sets of Dirac matrices
from their definitions \\
\texttt{imat.rs} & 316 & exact matrix arithmetic over the integers and
the Gaussian integers; exact rank \\
\texttt{verify.rs} & 582 & the machine checks of the matrices
(anticommutators, Lie algebra, group, notebooks, chirality) \\
\texttt{dirac.rs} & 272 & the Dirac equations of the two fields, checked
in momentum space \\
\texttt{pairs.rs} & 556 & which pairs of two identical particles exist;
their quantum numbers \\
\texttt{symmetry.rs} & 551 & internal symmetries of the zero-mode
sector \\
\texttt{kernel.rs} & 463 & exchange amplitudes of the Vorton-Scalar,
channel by channel \\
\texttt{gap.rs} & 414 & the gap equation of the hidden-axis singlet \\
\texttt{dimer.rs} & 171 & two particles in the Yukawa potential \\
\texttt{cjt\_toy.rs} & 201 & the method of Cornwall, Jackiw and
Tomboulis on a solvable model \\
\texttt{cjt\_pairing.rs} & 778 & effective potential of the paired
state; thermodynamics; landscape \\
\texttt{cjt\_normal.rs} & 807 & effective potential for the expectation
value of the Vorton-Scalar: matter, temperature \\
\texttt{cjt\_vacuum.rs} & 163 & the gap equation in the vacuum \\
\texttt{cjt\_cmd.rs} & 612 & the command \texttt{cjt} \\
\texttt{cosmo.rs} & 1108 & expansion history, the dark radiation of the
model, recombination \\
\texttt{boltzmann.rs} & 727 & linear perturbations: the
Einstein-Boltzmann equations \\
\texttt{spectra.rs} & 358 & angular power spectra and matter power
spectrum \\
\texttt{bessel.rs} & 222 & spherical Bessel functions for the
line-of-sight integrals \\
\texttt{forecast.rs} & 535 & mock likelihood, Fisher matrices, profile
fits \\
\texttt{fmat.rs} & 187 & small real matrices in double precision \\
\texttt{quad.rs} & 146 & Gauss-Legendre quadrature \\
\texttt{export.rs} & 136 & writing matrices and check reports \\
\texttt{main.rs} & 978 & the commands \\
\texttt{lib.rs} & 56 & list of the modules \\
\texttt{complex.rs}, \texttt{linalg.rs}, \texttt{roots.rs},
\texttt{io.rs} & 745 & complex numbers, small Hermitian eigenproblems,
root finding, CSV output (from NCSS) \\
\end{longtable}

\hypertarget{the-python-scripts-and-tools}{%
\section{5. The Python scripts and
tools}\label{the-python-scripts-and-tools}}

\begin{longtable}[]{@{}
  >{\raggedright\arraybackslash}p{(\columnwidth - 2\tabcolsep) * \real{0.4432}}
  >{\raggedright\arraybackslash}p{(\columnwidth - 2\tabcolsep) * \real{0.5568}}@{}}
\toprule\noalign{}
\begin{minipage}[b]{\linewidth}\raggedright
Script
\end{minipage} & \begin{minipage}[b]{\linewidth}\raggedright
Purpose
\end{minipage} \\
\midrule\noalign{}
\endhead
\bottomrule\noalign{}
\endlastfoot
\texttt{code/scripts/verify\_matrices\_independent.py} & independent
verification of the matrices: rebuilds them from the definitions with
exact rational arithmetic and repeats the essential checks (36 checks);
writes \texttt{data/matrix\_checks\_independent.txt} \\
\texttt{code/scripts/validate\_cmb\_against\_class.py} & compares the
spectra of the Rust program with reference spectra of the code CLASS;
writes \texttt{cmb\_validation.csv}, \texttt{pk\_validation.csv} \\
\texttt{code/scripts/make\_tables.py}, \texttt{make\_tables\_cosmo.py},
\texttt{make\_tables\_cjt.py} & turn the data files into the numbers and
tables of the documents (\texttt{provenance/generated/tables.md}) \\
\texttt{code/scripts/make\_figures.py} & draws all figures from the data
files (Sec. 7) \\
\texttt{code/scripts/assemble\_docs.py} & inserts the generated numbers,
tables and files into the document sources
\texttt{docs\_src/*.src.md} \\
\texttt{code/scripts/build\_docs.sh}, \texttt{check\_layout.py},
\texttt{doc\_header.tex} & Markdown to LaTeX to PDF for the supplements
and provenances, with a check of the layout \\
\texttt{code/scripts/build\_manuscripts.py}, \texttt{tex2md.py} & build
the two manuscripts from the common REVTeX source in
\texttt{docs\_src/manuscript/}; Markdown versions \\
\texttt{code/scripts/write\_symbol\_registry.py},
\texttt{check\_symbols.py}, \texttt{symbol\_concordance.py} & the
registry of symbols and its checks (provenance P1) \\
\texttt{tools/nb2txt.py}, \texttt{tools/extract\_nb\_matrices.py},
\texttt{tools/scan\_notebook\_symbols.py} & reading the notebooks
without Mathematica (provenance P3) \\
\end{longtable}

\hypertarget{from-a-data-file-to-a-number-in-a-document}{%
\section{6. From a data file to a number in a
document}\label{from-a-data-file-to-a-number-in-a-document}}

No number in the supplements, in ``dirac matrices'' or in the
manuscripts is typed by hand. The chain is

\begin{enumerate}
\def\labelenumi{\arabic{enumi}.}
\tightlist
\item
  a command of the Rust program writes a data file (Sec. 3);
\item
  \texttt{make\_tables.py} reads the data files and writes named
  fragments (a number, a table) into
  \texttt{provenance/generated/tables.md}; \texttt{spacetimes\ matrices}
  writes the displayed matrices into
  \texttt{provenance/generated/matrices.md};
\item
  the document sources contain placeholders (the name of a fragment in
  double curly braces, preceded by \texttt{GEN:}, \texttt{GENM:} or
  \texttt{ROWS:}); \texttt{assemble\_docs.py} and
  \texttt{build\_manuscripts.py} replace them.
\end{enumerate}

To find the origin of a number, search for it in
\texttt{provenance/generated/tables.md}; the name of the fragment leads
to the function of \texttt{make\_tables*.py} that made it, and that
function names the data file and the column. To change a number, the
program must be changed and run again; editing a document does not
survive the next build.

\hypertarget{from-a-data-file-to-a-figure}{%
\section{7. From a data file to a
figure}\label{from-a-data-file-to-a-figure}}

All figures are drawn by \texttt{code/scripts/make\_figures.py}, one
function per figure; each figure exists as PDF (for the documents) and
PNG.

\begin{longtable}[]{@{}
  >{\raggedright\arraybackslash}p{(\columnwidth - 4\tabcolsep) * \real{0.2553}}
  >{\raggedright\arraybackslash}p{(\columnwidth - 4\tabcolsep) * \real{0.3404}}
  >{\raggedright\arraybackslash}p{(\columnwidth - 4\tabcolsep) * \real{0.4043}}@{}}
\toprule\noalign{}
\begin{minipage}[b]{\linewidth}\raggedright
Figure (\texttt{figures/})
\end{minipage} & \begin{minipage}[b]{\linewidth}\raggedright
Data files
\end{minipage} & \begin{minipage}[b]{\linewidth}\raggedright
Shows
\end{minipage} \\
\midrule\noalign{}
\endhead
\bottomrule\noalign{}
\endlastfoot
\texttt{fig\_kernels} & \texttt{pair\_channel\_kernel.csv} & exchange
amplitudes in the pair channels \\
\texttt{fig\_gap} & \texttt{gap\_scan.csv} & the gap and the critical
temperature \\
\texttt{fig\_dimer} & \texttt{dimer.csv} & scattering length and binding
energy of two fermions \\
\texttt{fig\_cjt\_toy} & \texttt{cjt\_toy.csv} & the effective action on
a solvable model \\
\texttt{fig\_cjt\_thermo} & \texttt{cjt\_thermo.csv} & gap, entropy and
heat capacity against temperature \\
\texttt{fig\_cjt\_landscape} & \texttt{cjt\_landscape\_scan.csv} &
energy as a function of the order parameter \\
\texttt{fig\_cjt\_normal} & \texttt{cjt\_normal.csv} & fermion mass in
matter; self-bound matter \\
\texttt{fig\_cjt\_thermal} & \texttt{cjt\_thermal\_profile.csv} & the
fermion mass at finite temperature \\
\texttt{fig\_cjt\_vacuum} & \texttt{cjt\_vacuum.csv} & the critical
coupling for a condensate in the vacuum \\
\texttt{fig\_cjt\_vacpot} & \texttt{cjt\_vacuum\_potential.csv} & the
potential of the Vorton-Scalar in the vacuum \\
\texttt{fig\_cmb\_validation} & \texttt{cmb\_validation.csv} & spectra
of this work against CLASS \\
\texttt{fig\_step\_thermo} & \texttt{step\_thermo.csv} & the dark
radiation of the model through its transition \\
\texttt{fig\_dr\_profiles\_planck}, \texttt{fig\_dr\_profiles\_deep} &
\texttt{dr\_profiles.csv} & fits of the dark-radiation families to the
two mock data sets \\
\texttt{fig\_dr\_residuals} & \texttt{dr\_residuals.csv} & residual
spectra \\
\texttt{fig\_dr\_zt} & \texttt{dr\_zt\_scan.csv} & dependence on the
transition redshift \\
\texttt{fig\_sfdm\_transfer} & \texttt{sfdm\_transfer.csv},
\texttt{sfdm\_summary.csv}, \texttt{sfdm\_bounds.csv} & suppression of
the matter power spectrum by the sound of a superfluid \\
\texttt{fig\_dm\_windows} & \texttt{dm\_halo\_estimates.csv},
\texttt{dm\_selfinteraction.csv}, \texttt{sfdm\_bounds.csv} & the
fermion masses that a galactic superfluid needs and those that
observations allow \\
\end{longtable}

\hypertarget{how-the-code-is-tested}{%
\section{8. How the code is tested}\label{how-the-code-is-tested}}

\begin{longtable}[]{@{}
  >{\raggedright\arraybackslash}p{(\columnwidth - 6\tabcolsep) * \real{0.2621}}
  >{\raggedright\arraybackslash}p{(\columnwidth - 6\tabcolsep) * \real{0.4466}}
  >{\raggedright\arraybackslash}p{(\columnwidth - 6\tabcolsep) * \real{0.1165}}
  >{\raggedright\arraybackslash}p{(\columnwidth - 6\tabcolsep) * \real{0.1748}}@{}}
\toprule\noalign{}
\begin{minipage}[b]{\linewidth}\raggedright
Test
\end{minipage} & \begin{minipage}[b]{\linewidth}\raggedright
What it establishes
\end{minipage} & \begin{minipage}[b]{\linewidth}\raggedright
Result
\end{minipage} & \begin{minipage}[b]{\linewidth}\raggedright
Where
\end{minipage} \\
\midrule\noalign{}
\endhead
\bottomrule\noalign{}
\endlastfoot
106 machine checks of \texttt{spacetimes\ matrices} (96 exact, 10
numerical) & the properties of the two sets of Dirac matrices, including
identity with the matrices printed in the notebooks & all pass &
\texttt{data/matrix\_checks.txt} \\
36 checks of the independent Python program & the same essential
properties by other methods & all pass &
\texttt{data/matrix\_checks\_independent.txt} \\
63 unit tests (\texttt{cargo\ test\ -\/-release}) & each numerical
routine against a closed result or an independent method & all pass &
the \texttt{tests} modules of the source files \\
comparison with CLASS & the Einstein-Boltzmann solver against reference
spectra & differences in the tables of S3 Sec. 2 and of M App. D &
\texttt{data/cmb\_validation.csv}, \texttt{data/pk\_validation.csv} \\
consistency checks through the effective action & the gap equation, the
bound-state threshold and the Ginzburg-Landau coefficients by a second
route & agreement stated in S4 Sec. 3 &
\texttt{data/cjt\_pairing\_check.csv}, \texttt{cjt\_threshold.csv},
\texttt{cjt\_gl\_summary.csv} \\
9 checks of the symbols & uniqueness of the symbols of the project & all
pass & \texttt{data/symbol\_checks.txt} \\
layout checks & no line, table or formula wider than its box; no
undefined reference & reported at every build & output of
\texttt{build\_docs.sh}, \texttt{build\_manuscripts.py} \\
\end{longtable}

The code CLASS itself was not run in this project: the reference spectra
are those of the author's project NCSS, computed there with CLASS
v3.4.2. Every error that the tests found is recorded in the error log
(\texttt{provenance/P9\_error\_log}).

\end{document}
